\documentclass[a4paper,12pt]{article}
\usepackage[utf8]{inputenc}
\usepackage[T1]{fontenc}
\usepackage[portuguese]{babel}
\usepackage{geometry}
\usepackage[dvipsnames]{xcolor}
\usepackage{soul}
\usepackage{hyperref}
\usepackage{microtype}
\geometry{margin=2.5cm}

% \hl usa soul que permite quebra de linha dentro do highlight
\sethlcolor{yellow}

\title{Revista Militar -- Geopolítica e Futebol}
\author{Trabalho Prático 1 --- Scripting no PLN}
\date{2025/26}

\begin{document}

\maketitle

\begin{abstract}
As seguintes frases foram selecionadas através de um modelo de linguagem
baseado em \textit{n-grams} (trigramas) com \textit{scoring} por
log-probabilidade média e suavização de Laplace:
\begin{itemize}
  \item Estava dado o mote para uma guerra que iria durar quatro dias.
  \item Um português em Inglaterra a fazer publicidade a uma empresa Americana, prossegue assim, via futebol, a nossa vocação Euro-Atlântica.
  \item Segundo Jacques Attali vivemos na época do poder económico e a indústria constitui-se como a sua base única.
\end{itemize}
\end{abstract}

\section{Texto da Fonte}

A indústria do futebol movimenta hoje uma economia de 160 mil milhões de euros em todo o mundo. Apenas 22 países têm um PIB superior a este total.
 Num mundo em constante mutação há fenómenos que extravasam a sua essência e transpõem as suas fronteiras para áreas que, em princípio, lhe estariam vedadas. As grandes competições desportivas da actualidade constituem um exemplo deste fenómeno e são, enquanto tal, mais um produto da globalização. Não só os \hl{Jogos Olímpicos} e os \hl{Campeonatos do Mundo de futebol}, com um a “bater-nos à porta”, mas também um conjunto cada vez maior de competições desportivas (algumas de âmbito estritamente nacional, como seja a liga norte-americana de basquetebol, a NBA), atingem actualmente um vastíssimo público telespectador, dando a conhecer novas realidades de outros continentes e povos. Em praticamente todas as nações do planeta, centenas de milhões de indivíduos compartilham as imagens desta poderosa e crescente, por assim dizer, indústria do entretenimento, a qual age intensamente na cultura e na economia das nações.
 Em termos geopolíticos não estão muito distantes os anos em que os \hl{Jogos Olímpicos} pareciam o “campo de batalha” do mundo bipolar então existente. A “luta” pelas medalhas desenrolava-se de uma forma alucinante e todos os dias da competição o número de medalhas dos EUA e \hl{ex-URSS} subiam, acabando inevitavelmente por, entre os dois, conquistarem a maioria do total em “jogo”, ou seja, através do desporto acentuavam a sua condição de superpotências.
 Quanto ao outro grande pólo de congregação das nações a nível mundial, os \hl{Campeonatos do Mundo de Futebol}, assistiu-se à afirmação e expansão geográfica dos nacionalismos até à primeira metade do século XX, à mesma “\hl{guerra fria}” ESTE-OESTE,
 se bem que aqui uma das superpotências não tivesse um papel muito marcante e, após a queda do \hl{muro de Berlim}, à cogregação das novas nações resultante da diplomacia de uma das maiores organizações mundiais, a \hl{FIFA}.
Actualmente, quando pensamos no futebol, para além do desporto em si, estamos a pensar em economia, globalização, afirmação de nacionalidade, política, solidariedade, mercantilismo, violência, racismo e hooliganismo. Face a tal conjectura pode-se mesmo afirmar que o futebol já é mais que um simples desporto é um verdadeiro espelho da sociedade e com a organização do \hl{Campeonato Mundo de 2010} agendada para a \hl{África do Sul}, o futebol irá mesmo chegar aos “quatro cantos do mundo”.
 Dada a maior representatividade nas diversas competições nacionais e internacionais e pelo crescente nível de atenção que os media e mundo empresarial lhe dão e até pela realização do mundial de 2006 na \hl{Alemanha}, o presente artigo ir-se-á debruçar exclusivamente sobre o futebol como desporto de massas, preconizando uma ligação entre o mesmo e a \hl{Geopolítica}. Desta análise serão excluídos os factores como - a solidariedade, a violência, a corrupção e o hooliganismo - que, pela sua envolvente, carecem de outro tipo de estudo.
Por volta do ano de 2 500 a.c., o imperador chinês \hl{Cheng-Ti} incluiu no treino militar dos seus soldados uma espécie de jogo rude e violento, no qual dois grupos disputavam, com as mãos e os pés, a posse de um balão feito de couro e recheado de crina de cavalo.
 Mais tarde, em pleno século XII, surgiu, em \hl{Inglaterra}, uma prática desportiva na qual se usava uma bola e que se designada mob football, esta prática, dada a sua violência, foi banida por decreto real. Na segunda metade do século XVIII veio a ser retomada nos colégios ingleses, embora de forma mais moderada, onde a prática de jogos viris que exigiam maior empenho muscular, tomava o lugar das actividades mais nobres que implicavam destreza e equilíbrio (equitação, esgrima, arco e caça, entre outras).
 O futebol só surge na sua forma contemporânea, também em \hl{Inglaterra}, por volta dos finais do século XIX, em plena época da revolução industrial e do capitalismo.
Rapidamente os proprietários das fábricas compreenderam as vantagens que poderiam advir deste desporto que então despontava: forte união dos trabalhadores em torno da empresa e reconhecimento rápido da mesma. A partir deste momento o futebol terá quatro etapas de desenvolvimento conceptual e geo-histórico: a expansão, a institucionalização, a globalização e a industrialização.
 Após a \hl{“descoberta”} \hl{Britânica}, o futebol depressa “invadiu” a \hl{Europa} entre 1870 e 1890 essencialmente pelos portos e ferrovias. Note-se que os clubes mais antigos e
 representativos de cada país se encontram em cidades portuárias ou nas grandes capitais. Em \hl{França} o \hl{Le Havre} e o \hl{Marselha}, em \hl{Espanha} o \hl{Atlético de Bilbao} e o \hl{Barcelona} e no caso \hl{Português o Benfica}, o \hl{Sporting} e o \hl{Porto}. Do mesmo modo foram as cidades portuárias como - \hl{São Paulo}, \hl{Rio de Janeiro}, \hl{Buenos Aires} e \hl{Montevideu} - são as primeiras terem contacto com este novo desporto que despontava na \hl{América do Sul}.
 Graças ao conceito da livre e fácil representatividade, banqueiros, comerciantes e industriais estão na origem da maior parte dos clubes que se espalharam por toda a \hl{Europa} entre 1890 e 1910. É neste permeio que surge, em 1904, a \hl{FIFA}, à qual iremos dedicar especial atenção em área própria deste artigo. Quatro anos depois surge a primeira competição internacional administrada e organizada pela \hl{FIFA}, mas integrada nos \hl{Jogos Olímpicos de Londres} de 1908. Esta foi ganha pelos “inventores” do jogo, a \hl{Inglaterra}, que viria a repetir o triunfo nos \hl{Jogos Olímpicos de Estocolmo} em 1912.
 Após a \hl{primeira guerra mundial} a era cosmopolita do futebol transformar-se-á numa era de nacionalismos, passando o futebol a assumir a dimensão de uma nação.
 Cada país revê-se na sua selecção e surge, em 1930 no \hl{Uruguai}, o primeiro \hl{Campeonato do Mundo de Futebol} organizado em exclusivo pela \hl{FIFA}, perdendo assim o vínculo que o ligava aos \hl{Jogos Olímpicos} e ao factor amadorismo que ainda hoje, embora como um sofisma, está presente no designado “espírito” \hl{Olímpico}.
 Seguiram-se em 1934 na \hl{Itália} e em 1938 na \hl{França} os 2º e 3º \hl{Campeonatos do Mundo de Futebol}, que dão, na altura, um cariz mundial a este desporto.
 Com a \hl{segunda guerra mundial} vem a paragem dos \hl{Campeonatos Mundiais de Futebol}, os quais viriam a ser retomados em 1950 no \hl{Brasil}. Este já envolveu uma escala mais global com as selecções a serem qualificadas para o campeonato através de jogos de apuramento em vários continentes.
 O futebol institucionaliza-se e fazia agora parte, em definitivo, das relações internacionais entre nações.
 c. \hl{A Globalização}
 A partir do \hl{Campeonato do Mundo} na \hl{Suiça} em 1954, o primeiro a ser televisionado, inicia-se a chamada era da globalização. É também neste ano que ocorre a fundação da \hl{UEFA}, sobre a qual nos iremos debruçar mais à frente.
 No \hl{Mundial de} 1966 em \hl{Inglaterra}, de boa memória para \hl{Portugal}, dá-se então o primeiro “boom” na lobalização mediática do futebol com 400 milhões de telespectadores a assistirem em directo à final estimando-se que o campeonato, no seu conjunto, tenha sido visto por dois mil milhões. A partir de então conjugaram-se dois factores primordiais para o crescimento e mundialização deste desporto - avanço tecnológico dos meios de comunicação e as competições internacionais entre clubes que se tinham iniciado em
1955 - possibilitando que, de uma forma mais célere, as massas se identificassem com o mesmo.
 Em 1998 no \hl{Mundial} em \hl{França} a audiência acumulada atingiu um total e 33,4 mil milhões de telespectadores e em 2002 na \hl{Coreia do Sul} e \hl{Japão} 28,8 mil milhões, como se pode verificar no \hl{Anexo A.}
 No “campo” da geografia, a globalização atingiu um marco histórico com o \hl{Mundial de}
1994 nos EUA onde, pela primeira vez, um campeonato decorria fora do continente
 \hl{Europeu} ou da \hl{América Latina}. Esse marco foi continuado em 2002 na \hl{Coreia do Sul} e \hl{Japão} e culminará em 2010 na \hl{África do Sul}.
 Paralelamente à globalização e a partir de meados da década de 1980 começa a nascer uma nova vertente do futebol, a industrialização. O crescente número de canais televisivos vieram, para além do número de telespectadores, aumentar a concorrência entre os mesmos e por consequência fizeram disparar o valor dos direitos televisivos dos jogos. A título de exemplo refira-se que entre 1985 e 1996 os direitos de transmissão de jogos da liga \hl{Inglesa} passaram dos 1,8 para os 250 milhões de dólares anuais.
 \hl{Começam} as transferências sonantes, em 1984 \hl{Diego} \hl{Maradona} transfere-se do \hl{Barcelona} para o \hl{Nápoles} (\hl{Itália}) por 11 milhões e 850 mil euros, em 2000 \hl{Luís Figo} transfere-se do \hl{Barcelona} para o \hl{Real Madrid} (\hl{Espanha}) por 63 milhões de euros e a escalada atinge o seu ponto máximo em 2001 com a transferência de \hl{Zinedine Zidane da Juventus} (\hl{Itália}) para o \hl{Real Madrid} por 75 milhões de euros a que se seguiram o \hl{Brasileiro} \hl{Ronaldo} e o \hl{Inglês} \hl{Beckman}.
 Na última década surge ainda um facto novo, as figuras do futebol começam a “vender” a sua “imagem” ao mundo empresarial que vê nos primeiros uma forma de expandir as suas vendas e de alcançar novos mercados. Os jogadores de “1ª linha” recebem hoje em dia mais com os direitos da sua “imagem” do que propriamente com os seus vencimentos.
Um exemplo é \hl{Ronaldo} que dos 26 milhões de euros/ano que recebe, dois terços são provenientes de contratos publicitários e de cedência de imagem a empresas de cariz mundial como a \hl{Siemens}, a \hl{Nike} e a \hl{Audi}. O \hl{“nosso}” \hl{José Mourinho}, treinador do \hl{Chelsea de Inglaterra}, foi considerado como a “marca comercial” mais cobiçada no \hl{Reino Unido} e assinou, em 2005, um contrato de 21 milhões de euros para aparecer num spot publicitário de um cartão de crédito \hl{Norte-Americano}. Um português em \hl{Inglaterra} a fazer publicidade a uma empresa \hl{Americana}, prossegue assim, via futebol, a nossa vocação \hl{Euro-Atlântica}.
 Os clubes face à \hl{Lei Bosman}, que liberaliza as transferências de jogadores e origina o aumento dos salários dos jogadores, e às “loucuras” dos seus dirigentes vêm-se endividados e a caminho da falência e \hl{Portugal} não foge à regra. Os governos começam a aperceber-se dos milhões de euros que o futebol movimenta e exigem as devidas obrigações fiscais aos clubes. Estes compreendem então que têm que ser geridos com rigor e por referenciais dos modelos empresariais, sendo a última novidade do futebol a criação das \hl{Sociedades Anónimas Desportivas} (SAD´s), algumas cotadas na bolsa.
 Definitivamente o futebol de hoje é uma “indústria”.
 Neste capítulo iremos destacar duas associações que, pela sua actividade extra futebol, alargaram os horizontes deste desporto vencendo barreiras que a política por si ainda não resolveu. Como visto no capítulo anterior a \hl{FIFA} tem um destaque especial no âmbito do futebol e de certa forma poder-se-á considerar que é uma organização mais global que a própria ONU. Actualmente a \hl{FIFA} tem 2058 filiados (incluindo a \hl{Palestina}), ou seja, mais que os 1919 da ONU. Já a \hl{UEFA}, mais regional, teve não só um papel importante no desenvolvimento mediático do futebol, como também nas relações internacionais \hl{Europeias} durante o período da \hl{“Guerra-Fria}”, altura em que iniciou a sua actividade.
 Com o acto da assinatura, pelos representantes dos sete países fundadores (\hl{França}, \hl{Bélgica}, \hl{Dinamarca}, \hl{Holanda}, \hl{Espanha}, \hl{Suécia} e \hl{Suiça}), em 21Mai esta associação vai influenciar a história e o rumo do futebol até à actualidade. Em 1905 juntam-se-lhe as associações de futebol da \hl{Inglaterra}, \hl{Escócia}, \hl{País de Gales}, \hl{Irlanda}, \hl{Alemanha}, \hl{Áustria}, \hl{Itália} e \hl{Hungria}. No entanto no início de 1906 a \hl{FIFA} era apenas uma associação \hl{Europeia}, só mais tarde se dá a entrada dos primeiros membros não europeus - \hl{África do Sul} (1909), \hl{Argentina} e \hl{Chile} (1912) e \hl{Estados-Unidos} (1913) 10 - dando início à sua expansão mundial. No fim da \hl{primeira guerra mundial} a \hl{FIFA} contava com 20 associações filiadas.
 Em 1930 e pela “mão” do seu presidente \hl{Jules Rimet} (eleito em 1921) surge o primeiro acto de diplomacia praticado pela \hl{FIFA}. No sentido de obter a maior participação possível no \hl{Mundial do Uruguai} e dada a crise que se vivia, a \hl{FIFA} comprometeu-se a assumir os custos das passagens e alojamento dos participantes.
 Mesmo com esta premissa só quatro países europeus vieram a participar no evento \hl{França}, \hl{Bélgica}, \hl{Jugoslávia} e \hl{Roménia} - provocando a ira do \hl{Uruguai} que decidiu não defender o seu título no campeonato seguinte em \hl{Itália}.
 No fim da era \hl{Jules Rimet} e na década seguinte, a geografia mundial do futebol ficou definida com a fundação de cinco das seis confederações que constituem a \hl{FIFA}, à CSF, já existente e fundada em 1916, juntam-se a \hl{UEFA} e a AFC em 1954, a CAF em 1957, a \hl{CONCACAF} em 1961 e a OFC em 1966. A distribuição geográfica ao nível mundial destas confederações encontra-se materializada no \hl{Apêndice 1}. Num olhar mais atento, essa distribuição praticamente que coincide com a divisão preconizada pelo pensador holandês \hl{Nicholas Spykman} após a \hl{II guerra mundial}. Na sua visão geopolítica global, \hl{Spykman}, preconizava a divisão do mundo em cinco grandes ilhas - continentes
(\hl{Apêndice 2):} \hl{América do Norte} e \hl{Eurásia}, os mais importantes, \hl{Austrália}, \hl{África} e \hl{América do Sul}. As principais diferenças residem no facto da \hl{Eurásia} de \hl{Spykman} ter sido desdobrada, pela \hl{FIFA}, em \hl{Europa} e \hl{Ásia} e também na \hl{“substituição}” da \hl{América do Norte} pela \hl{América do Sul} em termos de importância.
 Quando em 1974 \hl{João Havelange} foi eleito presidente da \hl{FIFA} o futebol passou a ser visto não só como uma competição, mas também como um veículo de desenvolvimento
 técnico que devia procurar novas formas de financiamento. Em pouco tempo a associação, sedeada em \hl{Zurique}, transformou-se numa “empresa” dinâmica à “conquista” do mundo económico e político.
 As ondas criadas com as convulsões políticas, particularmente em \hl{África}, onde antigas colónias se tinham constituído em países, depressa chamou a atenção da \hl{FIFA} e no \hl{Campeonato do Mundo de 1982} em \hl{Espanha} o número de participantes foi alargado para
24, permitindo a entrada de mais selecções da \hl{África}, \hl{Ásia} e \hl{América do Norte} e \hl{Central} que, até então, só tinham direito ao envio de uma selecção por confederação. O alargamento foi considerado um sucesso e uma afirmação de que as ideias e políticas então estabelecidas estavam correctas. No \hl{Mundial de 1998} em \hl{França} o número de participantes subiu para 32 permitindo uma maior participação de todas as confederações que constituem a \hl{FIFA}.
 Em termos políticos esta associação, que preconiza o principio da universalidade e ao qual está comprometida, tornou-se então num eixo de diplomacia conseguindo em 1975 a adesão da \hl{China} sem que para tal tivesse que excluir a \hl{Formosa} e em 1991 consegue que a \hl{Coreia do Sul} e \hl{Coreia do Norte} enviem uma selecção conjunta ao \hl{Mundial de juniores de Portugal}. O caso mais paradigmático desta diplomacia política consistiu na inclusão de \hl{Israel} na \hl{UEFA} face à sua situação no médio oriente. Como culminar destas acções e no intuito de reconciliar o \hl{Japão} e a \hl{Coreia do Sul}, cujas relações são historicamente difíceis, e ao mesmo tempo facilitar a \hl{reunificação da Coreia}, atribui a organização do \hl{Campeonato do Mundo de 2002} conjuntamente ao \hl{Japão} e à \hl{Coreia do Sul}.
 Para melhor caracterizar a \hl{FIFA} diplomática, fica aqui uma declaração de \hl{João Havelange} antes do \hl{Campeonato do Mundo de 1998}: “Há um projecto que não se realizou mas que eu espero vir a concretizar. Seria um desafio entre as selecções da \hl{Palestina} e de \hl{Israel}.”
 Em 1954 surge, na \hl{Europa}, a associação que viria a revolucionar o futebol ao nível regional e de clubes. Estamos em plena \hl{“Guerra-Fria}”, onde predomina o “fantasma” nuclear, onde cruzar fronteiras na designada cortina de ferro é extremamente complexo e onde a \hl{OTAN} e o \hl{Pacto de Varsóvia} intensificam as suas actividades de espionagem. É neste contexto que a \hl{UEFA}, inicialmente com 30 federações nacionais filiadas, cria as competições internacionais de clubes e consegue, de forma pacífica, “eliminar” as fronteiras com o \hl{Bloco de Leste} via futebol.
 A \hl{Taça dos Campeões}, a \hl{Taça das Feiras} e a \hl{Taça UEFA} vieram possibilitar a existência de uma \hl{Europa} única, sonho político ainda hoje por atingir. Todo o continente \hl{Europeu} respeita as mesmas regras do futebol e as poucas imagens e notícias disponíveis do \hl{Bloco de Leste} surgem através do mesmo.
 Em 1960 inicia-se o \hl{Campeonato da Europa} de selecções sob a égide da \hl{UEFA} e podemos ver uma selecção da \hl{URSS} com jogadores da \hl{Rússia}, \hl{Ucrânia} e \hl{Lituânia} representando o
 mesmo símbolo e cantando o mesmo hino. Em 1991 podemos ver o \hl{Estrela Vermelha de Belgrado} vencer a \hl{Taça dos Clubes Campeões Europeus} com jogadores de origem \hl{Sérvia}, \hl{Croata} e \hl{Bósnia} a levantar, em simultâneo, o troféu.
 Com a implosão geopolítica da \hl{URSS} e da \hl{Jugoslávia} rapidamente o que era já não o é.
Após a declaração de independência, em 1990, a \hl{Lituânia} retirou a sua equipa de futebol do campeonato soviético e outros países se seguiram, como consequência passaram a haver novos campeões nacionais. Um dos primeiros passos para o reconhecimento de um país como que passou a ser a sua entrada na \hl{UEFA}, que cresceu a tal ponto de hoje contar com 5213 federações nacionais filiadas, ou seja mais do dobro dos 25 países que constituem a \hl{União Europeia} após o último alargamento.
 As organizações e competições internacionais de futebol, atrás referidas, dão, certamente, uma imagem mais real da mundialização do futebol através da presença, nas mesmas, de todos os continentes, países e regiões. Um dos resultados práticos dessa mundialização é o facto dos \hl{“astros}” do futebol serem mais conhecidos que os próprios chefes de estado. Estes últimos já se aperceberam deste fenómeno e começam cada vez mais a aproximar-se dos ícones futebolísticos dos respectivos países. \hl{Em Portugal} a selecção nacional, após o desempenho no \hl{EURO 2004}, foi recebida em \hl{Belém} como sinal de reconhecimento de feitos nacionais, na \hl{Alemanha} é \hl{Franz Beckenbauer} que está a liderar a organização do \hl{Mundial de 2006 naquele país}, na \hl{França} \hl{Michel Platini} acompanhou em 1997 o presidente \hl{Jacques Chirac} numa visita a vários países da \hl{MERCOSUL16}, na qual foram restabelecidas as relações diplomáticas entre a \hl{França} e a \hl{América Latina}.
 Como corolário desta mundialização do futebol e consequente projecção da imagem de um país, o pequeno \hl{Marcunis} foi encontrado, 15 dias após o tsumani que varreu a província de \hl{Ache} na \hl{Indonésia}, com uma camisola da selecção \hl{Portuguesa} que tinha desde o \hl{Mundial de 2002} na \hl{Coreia}.
 Não só o povo se exprime através do futebol, o mundo político também o faz e neste subcapítulo iremos evidenciar determinados acontecimentos no mundo do futebol que tiveram ligação directa à política ao longo do último século.
 Em 1934 a \hl{Itália} fascista de \hl{Mussolini} utiliza o campeonato como meio de propaganda do regime e de aceitação internacional, chegando ao ponto de designar um estádio em \hl{Roma} como “O \hl{Estádio do Partido Fascista}”. Anos mais tarde o \hl{General Franco} impede a \hl{Espanha} de jogar contra a \hl{Rússia} no \hl{Campeonato da Europa} de 1960. Em 1978, por altura do \hl{Mundial} na \hl{Argentina}, organiza-se um comité de boicote ao evento para alertar a comunidade internacional para os malefícios da junta militar de \hl{Jorge Videla}, no
 entanto a selecção da “casa” viria a ganhar o mesmo e a fazer esquecer o assunto, não sem antes ter acontecido o \hl{“famoso”18 Argentina-Perú}.
 Não só as selecções mas também os clubes foram utilizados como “ferramentas” da política, basta recordar as conquistas do \hl{Real Madrid} na década de 50, para se perceber qual era o clube do \hl{Franquismo}. Esta situação inflamava as identidades regionais da \hl{Catalunha} e do \hl{País Basco} que eram exteriorizadas através dos jogos dos seus clubes mais representativos, o \hl{Barcelona} e o \hl{Atlético de Bilbau} respectivamente. Ainda hoje se fazem sentir estas clivagens nos jogos disputados entre o \hl{Barcelona} e o \hl{Real Madrid}. Já o \hl{Atlético de Bilbau} continua a só admitir na sua equipa jogadores de origem basca.
 A propaganda via futebol foi também uma realidade portuguesa com o \hl{Estado Novo de Salazar} e as conquistas internacionais do \hl{Benfica} e da selecção nacional na década de 60.
 Ainda na década de 60 começam a fazer-se sentir os ecos de \hl{África}, quando, na fase de apuramento para o \hl{Mundial de} 1966, quinze países \hl{Africanos} decidem retirar as respectivas selecções em protesto pela sub-representatividade daquele continente na competição.
 Em 1969 dá-se, porventura, o caso mais representativo da mistura de futebol, política e nacionalismo. As \hl{Honduras} e \hl{El Salvador} tiveram de se enfrentar nos jogos de qualificação para o \hl{Campeonato do Mundo de 1970}. Tendo perdido por 1 a 0 no desafio da 1ª mão, \hl{El Salvador} acabou por vencer a eliminatória na 2ª mão num ambiente de tal forma hostil que a equipa das \hl{Honduras} teve de ser acompanhada ao estádio em viaturas blindadas. Face a esta situação e à eliminação, milícias armadas expulsaram os camponeses salvadorenhos instalados nas \hl{Honduras}. Estava dado o mote para uma guerra que iria durar quatro dias. Neste caso o futebol forneceu a centelha da guerra mas o barril da pólvora já lá estava.
 Em plena \hl{“Guerra-Fria}” há clubes do leste europeu que nos ficam na memória por representarem e pertencem aos respectivos regimes, estamos a recordar, a título de exemplo, do \hl{Steaua de Bucareste}, da \hl{Roménia} de \hl{Nicolau Ceausescu}, que ganhou a \hl{Taça dos Campeões Europeus de 1986}.
 Em 1990 surgem confrontos entre adeptos do \hl{Dínamo de Zagreb} (croata) e o \hl{Partizan} de \hl{Belgrado} (sérvio) dos quais resultaram 60 feridos e em \hl{Split} (\hl{Croácia}), no decurso de um desafio entre o \hl{Hadjuk} e o \hl{Estrela Vermelha de Belgrado}, a bandeira jugoslava foi simbolicamente queimada pelos croatas. Estas acções eram o prenúncio do que viria a acontecer no mês de \hl{Julho} do ano seguinte nos \hl{Balcãs}. O futebol poderia ter constituído um bom revelador das tensões existentes entre as diferentes repúblicas que formavam a \hl{Jugoslávia}. Considerada responsável pela guerra, tanto na \hl{Croácia} como na \hl{Bósnia}, a equipa jugoslava (\hl{Sérvia e Montenegro}, no futuro só \hl{Sérvia} devido ao recente referendo no \hl{Montenegro}) foi excluída do \hl{Campeonato da Europa} de 1992 para a qual tinha sido uma das oito equipas apuradas. Para a comunidade internacional, foi um meio de agir simbolicamente, de punir \hl{Belgrado} sem incorrer em qualquer risco militar.
 A \hl{África do Sul}, um dos primeiros países a aderir à \hl{FIFA}, foi, durante dezenas de anos, afastada das competições internacionais devido ao seu regime de apartheid. Como sinal de reconhecimento da sua mudança política acolhe em 1996, pela primeira vez, a \hl{Taça Africana das Nações}.
 Como se pode ver pelos exemplos aqui descritos o futebol, para além da sua vertente desportiva, acabou ao longo dos anos por ser um dos meios da política, quer pelas afirmações nacionalistas, quer pelos boicotes e sanções que ocorreram. Mas também há casos de ascensão ao poder via futebol e vejamos aqui um exemplo. Em 1986 \hl{Silvio Berlusconi} “pega” num \hl{AC Milan} à beira da falência e em poucos anos transformou-o num clube de sucesso. O sucesso do clube projectou-se de tal modo na imagem de \hl{Berlusconi} que, em 1994, forma o partido \hl{Forza Itália} e foi o \hl{Chefe de Governo} italiano até às últimas eleições que, a muito custo, reconheceu ter perdido para \hl{Romano Prodi}.
 Segundo \hl{Jacques Attali} vivemos na época do poder económico e a indústria constitui-se como a sua base única. Na sua visão prospectiva considera que o mundo evoluirá para uma sociedade hiper industrializada dominada por dois espaços - o espaço \hl{Europeu} e o espaço do \hl{Pacífico} (\hl{Japão} e EUA) - em que as potências económicas se substituirão às potências militares. O futebol parece ter interpretado à letra a esta visão geopolítica.
Depois da “aventura” americana do \hl{Mundial de 1994}, a \hl{FIFA “}levou” o \hl{Mundial de 2002} para a \hl{Coreia do Sul} e \hl{Japão}, não só pelos motivos políticos referidos anteriormente mas também, porventura em primeira prioridade, por motivos económicos.
 Com base nos conceitos de \hl{Attali} pode-se considerar que o coração económico do futebol se encontra essencialmente na \hl{Europa} e tenta atingir a \hl{Ásia}, depois de ter passado nos EUA. Após o \hl{Mundial de 2002} os “clubes marca” europeus parecem ter assumido em definitivo que o \hl{Oriente} é o local privilegiado para a sua expansão. O \hl{Real Madrid}, um dos clubes mais mediáticos dos últimos anos, apostou na contratação de “galáticos” facilmente reconhecidos na \hl{Ásia} como - \hl{Beckman}, \hl{Ronaldo}, \hl{Zidane} e \hl{Luís Figo} - e nas duas últimas temporadas em vez dos tradicionais estágios de inicio de época longe das multidões, “passeia” os seus “galáticos” por terras do \hl{Oriente} acompanhado pelos contratos milionários dos jogos de exibição e pelos milhares de fãs sempre prontos a comprar um “souvenir” da equipa. O resultado desportivo na última temporada é o que se conhece, acabou o campeonato com uma derrota, despediu dois treinadores e correu o risco de não atingir a milionária \hl{Liga dos Campeões}. Mas será que o fracasso desportivo importa face aos 300 milhões de euros que factura por ano. Outro exemplo é o do \hl{Manchester United} que abriu \hl{“Reds Cafés}” em \hl{Singapura} e na \hl{Malásia} onde os adeptos orientais podem ver os “seus heróis”, em directo, nos jogos da \hl{Liga Inglesa} e que há dois anos que perde a \hl{Liga Inglesa} para o \hl{Chelsea de Abramovic} e \hl{Mourinho}.
 Na periferia temos países como o \hl{Brasil}, a \hl{Argentina} e os países \hl{Africanos}, que possuem a “matéria-prima” mas não têm condições para a manter, tendo de a “exportar” para o
 coração. O número de jogadores da periferia a actuar no coração tem vindo a aumentar e assiste-se à “fuga” dos talentos da \hl{América do Sul} e de \hl{África} para o coração. De acordo com os cálculos da \hl{Confederação Brasileira de Futebol} há hoje em dia cinco mil brasileiros a actuar profissionalmente em clubes estrangeiros, incluindo a \hl{China}, \hl{Líbano}, \hl{Vietname} e até as \hl{Ilhas Faroe}.
 O caso mais paradigmático do poder económico do coração \hl{Asiático} estará porventura para acontecer. O primeiro-ministro tailandês propõe-se comprar 30\% do \hl{Liverpool}
(Clube da \hl{FA Cup Inglesa}) admitindo utilizar não só fundos privados como do próprio estado. Considera que tal negócio será um investimento uma vez que, para além de desenvolver o futebol tailandês, permitirá exportar uma imagem positiva do país através da publicidade daquele clube.
 O futebol apresenta ainda uma particularidade que contraria a lógica estratégica mundial: a superpotência americana não é dominante.
 Por pouco que a mundialização não se resume à americanização do mundo. Se pensamos na bolsa, vem-nos à cabeça \hl{Wall Street}, se pensamos em informação, surge-nos a CNN, se pensamos em cinema é com \hl{Hollywood} e as novas tecnologias estão no \hl{Silicon Valley}.
No entanto quando pensamos no futebol é preciso ser um especialista para saber qual a melhor equipa ou jogador da \hl{“Major League Soccer}” nos EUA.
 Apesar dos esforços desenvolvidos pela \hl{FIFA} e pela organização do \hl{Mundial de 1994}, o futebol masculino não se conseguiu impor aos desportos nacionais americanos como o basebol e o futebol americano, estes praticamente circunscritos àquele país. Já o futebol feminino é um exemplo de afirmação interna e internacional, no entanto a sua dimensão mundial ainda está muito atrasada.
 As superpotências do futebol, se considerarmos as vitórias dos últimos 40 anos dos campeonatos mundiais, são o \hl{Brasil}, a \hl{Alemanha}, a \hl{Itália}, a \hl{Argentina}, a \hl{França} e a \hl{Inglaterra}, logo o futebol representa um mundo multipolar que contraria a hegemonia americana, que vai da economia à tecnologia, passando pelo entretenimento. Neste mundo multipolar, no qual os EUA não são um pólo, os países \hl{Africanos} começam a tornar-se potências emergentes como foi o caso dos \hl{Camarões} e \hl{Senegal}, que chegaram aos quartos de final dos últimos campeonatos mundiais e a \hl{Ásia} começa a despontar com a \hl{Coreia do Sul} que atingiu as meias-finais em 2002.
Como se pôde verificar ao longo deste trabalho o futebol, para além do jogo em si, apresenta fortes ligações à economia e à política. Face a tal ingerência há factores da geopolítica que são afectados directamente pelo futebol como sejam o factor \hl{Humano} de
 cariz etnográfico e o factor Estruturas na sua vertente económica. Há ainda um terceiro factor que, por ocasião dos grandes eventos, assume particular relevância como seja o factor \hl{Circulação}. Este capítulo procura analisar a forma como estes três factores são influenciados pelo futebol, dando como exemplo, sempre que possível, a realidade \hl{Portuguesa} e a recente organização do \hl{Campeonato da Europa} de 2004.
 É neste factor que mais se faz sentir a “força” do futebol. Segundo dados retirados do site oficial da \hl{FIFA} referentes a 2000, praticam este desporto a nível mundial 242 milhões de pessoas correspondendo a 4,1\% da população, dentro destes incluem-se 52 milhões de europeus que correspondem a 6,7\% da população na \hl{Europa}. \hl{Em Portugal} estima-se que esse número chegue aos 291 mil. Perante estes factos é fácil compreender que todos os dias as televisões dêem noticias do futebol e que diariamente se publiquem jornais desportivos em que a maior parte das noticias são relativas ao futebol, cimentando, por assim dizer, a cultura do futebol.
 A popularidade à escala mundial e foco dos media, fazem com que o comportamento desportivo adquira um significado social, em especial quando estão em causa as selecções nacionais, formando uma coesão, sensação de poder e identidade nacional em torno da mesma. São múltiplos os casos de identificação nacional com a respectiva selecção, mas há um que, pela sua importância na altura, se revelou particularmente crítico. Por ocasião do \hl{Campeonato Europeu} de 1996 na \hl{Inglaterra} defrontaram-se, nas meias-finais, a equipa da casa e a \hl{Alemanha}, este jogo como que foi uma reedição dos combates travados ao longo da segunda grande guerra, com os média a apelar ao nacionalismo e a “espicaçar” a vontade popular, com os governantes preocupados com a segurança e com as autoridades a apelar à calma.
 \hl{Em Portugal} desde a revolução do \hl{25 de Abril de 1974} que não se via uma tão grande identificação da população com os símbolos nacionais - bandeira e hino nacional - como a que se viveu em torno da \hl{Equipa de Todos Nós} no último \hl{EURO 2004}. Foram as bandeiras penduradas nas janelas, foi o cantar do hino nacional a “plenos pulmões” no início de cada jogo, foram as manifestações populares em todo o país após as vitórias alcançadas e foi a presença dos mais altos dirigentes nacionais nos jogos da selecção portuguesa. A esta união não faltou também um relembrar da história e um apelo ao sentimento nacional quando o treinador \hl{Scolari} referiu que o jogo contra a \hl{Espanha} era de \hl{“Matar} ou morrer”.
 \hl{Veremos} o que sucederá no próximo mês de \hl{Junho} a reedição do \hl{EURO 2004} ou a frustração da crise económica.
 A terminar a análise deste factor cita-se o relatório síntese da avaliação do impacto do \hl{EURO 2004},“O projecto cultural do \hl{EURO2004} atingiu os objectivos inicialmente propostos, nomeadamente…, a promoção de orgulho cívico e auto-estima nacionais, a promoção de uma imagem positiva do país nos domínios imanentes à cultura no contexto
 da \hl{identidade europeia de Portugal}, …” 26
 Cada vez mais os governos vêm na organização dos eventos internacionais de futebol uma oportunidade para impulsionar as suas políticas económicas. O argumento usado pelo governo, aquando da decisão de organizar o \hl{Campeonato Europeu} de 2004 \hl{em Portugal}, sempre apontou para o facto de que os benefícios iriam compensar os custos. O governo alemão face à organização do \hl{Mundial de 2006}, investiu 6,5 mil milhões de euros em infra-estruturas: modernização dos estádios, construção e ampliação de estradas, sistemas electrónicos de monitorização de tráfego. Em contrapartida, conta com um crescimento de 8 mil milhões de euros no \hl{Produto Interno Bruto} (PIB) do país, entre 2003 e 2010, como consequência dos impulsos resultantes da realização do campeonato.
 As estruturas empresariais rapidamente tentam associar-se a este tipo de eventos tornando-se patrocinadores dos mesmos e procurando a “imagem” dos jogadores mais mediáticos para promover os seus serviços e produtos. É a \hl{“guerra}” para ver quem patrocina o quê.
 As televisões esforçam-se para obter o exclusivo dos direitos de transmissão. Analisando os dados da \hl{FIFA} verifica-se que a cada \hl{Campeonato do Mundo} os direitos televisivos como que se multiplicam. O \hl{Mundial} na \hl{Alemanha} em 2006 irá render a esta associação 1,5 mil milhões de \hl{Francos} \hl{Suíços}, só nos direitos televisivos adquiridos pelo \hl{KirchSport Group}, como se pode comprovar no \hl{Anexo A. Já} a \hl{UEFA}, mais comedida, apresentou receitas na ordem dos 556 milhões de euros em direitos televisivos decorrentes da \hl{Liga dos Campeões} na época de 2003/0428.
 A avaliação do impacto económico do \hl{EURO 2004} concluiu que o mesmo “… foi um evento de sucesso, …, na minimização de eventuais impactos negativos, em especial financeiros, do pós-evento”. A um investimento pré-evento de 964 milhões de euros em novos estádios, acessibilidades relacionadas e outros investimentos, contrapôs-se um valor acrescentado de 694 milhões de euros associado ao programa de investimentos e um valor acrescentado de 81 milhões de euros associado ao turismo.
 Por ocasião dos grandes eventos as organizações preocupam-se não só com os jogos em si como também com as acessibilidades para “chegar” aos mesmos. Por ocasião do \hl{EURO}
2004 foram criadas uma série de condições como novos acessos, alargamento de linhas de metro, autocarros e comboios especiais que custaram ao \hl{Estado} 113 milhões de euros.
 Das condições criadas algumas foram temporárias, no entanto as alterações de fundo não desapareceram com o evento, ficando aqui alguns exemplos representativos:
\hl{Remodelação} do nó de \hl{Francos} da \hl{Via de Cintura Interna} no \hl{Porto}, a nova ponte da
 \hl{Portela} em \hl{Coimbra} e a ligação da circular urbana de \hl{Guimarães} à variante de \hl{Fafe}. No que respeita aos sistemas de transportes públicos “ficou” a electrificação da ligação ferroviária para o \hl{Algarve} e a nova estação de \hl{Faro}, o \hl{Porto} tem duas novas estações de metro: \hl{Francos} a servir o \hl{Estádio do Bessa} e a estação do \hl{Estádio do Dragão}.
 Em cada competição internacional, certamente que nos países organizadores irão surgir planos especiais consoante as situações, mas também é certo que há “obras” que irão perdurar no tempo.
Desde o fim do século XIX, altura em que despontou para o mundo, que o futebol cresceu e sofreu várias transformações conceptuais não havendo dúvidas que hoje, em termos geográficos, atingiu a globalização. Também não existem dúvidas que, como descrito na frase de abertura deste trabalho, o futebol é uma indústria à volta da qual gravitam milhões de euros.
 O maior veículo de expansão deste desporto foi e é uma organização que conta mais membros do que qualquer organização política que conhecemos, a \hl{FIFA}, e à qual qualquer novo \hl{Estado-Nação} parece interessado em aderir. O século XX viu eclodir, por motivos vários, uma série desses \hl{Estados} que provocaram uma alteração \hl{Geopolítica} mundial e, como consequência, um aumento de equipas nacionais. Por entre as primeiras manifestações de vontade dos novos \hl{Estados} independentes figurava o pedido de adesão às organizações que regem o futebol, como se a definição de \hl{Estado} para além dos três elementos tradicionais: um território, uma população, um governo, necessitasse agora de um quarto: uma equipa nacional de futebol.
 Nos últimos 50 anos o futebol não só se transformou no maior desporto mundial, mas também cresceu noutras áreas da sociedade como a economia e a política interpretando da melhor forma as visões de geopolíticos consagrados. Com uma divisão multipolar do mundo, o futebol mais do que qualquer factor envolve regiões, pessoas e nações, fazendo parte da sua cultura.
 Com aproximadamente 240 milhões de jogadores activos e biliões de telespectadores, o futebol constitui uma substancial oportunidade industrial abrindo novos mercados para si e para o resto do mundo dos negócios. Os clubes estão cotados em bolsa, as empresas apostam na imagem dos \hl{“astros} da bola” para aumentar as suas vendas e os media concorrem ferozmente pelo maior “share” proporcionado pelas transmissões televisivas dos jogos de futebol.
 O futebol já não é só um desporto de massas, no entanto, não nos podemos esquecer que da “invenção” chinesa até aos dias de hoje, embora tenha mudado, não perdeu um carácter de simulação da guerra entre dois grupos rivais. De fato, o “pontapé na bola”
(foot-ball) carrega consigo a metáfora da contenda entre dois “exércitos” iniciada há 45 séculos atrás, a começar pelo próprio local em que se pratica: o campo (de batalha).
 Apesar de tudo continua ainda a ser nesse campo que se empregam as tácticas e estratégias do jogo, que se defende e ataca e que se utilizam as “armas” que levam às vitórias e às derrotas.

\bibliographystyle{plain}
\begin{thebibliography}{1}
  \bibitem{revista_militar}
  Geopolítica e o Desporto de Massas --- O Futebol.
  Disponível em: \url{https://www.revistamilitar.pt/artigopdf/90}.
  Acedido em março de 2026.
\end{thebibliography}

\end{document}
